%% kapa: exam=1 round=34 year=2023 period=2 form=A subject=appraisal_law subject_ko=감정평가관계법규 seq=32 number=32
%% source: Q-Net 감정평가사 기출문제 (cst00309) artlSeq=5211789 file=raw/1008/5211789/2023년도 제34회 감정평가사 1차시험 2교시.pdf
%% conversion: kordoc 4.24.51 -> tools/gen1.py (auto); review=proofread-sonnet; flags=none
\begin{question}{32}
부동산등기법령상 표시에 관한 등기에 관한 설명으로 옳지 않은 것은?
\begin{choices}
  \item 등기관은 토지 등기기록의 표제부에 등기목적을 기록하여야 한다.
  \item 토지의 분할이 있는 경우 그 토지 소유권의 등기명의인은 그 사실이 있는 때부터 1개월 이내에 그 등기를 신청하여야 한다.
  \item 구분건물로서 그 대지권의 변경이 있는 경우 구분건물의 소유권의 등기명의인은 1동의 건물에 속하는 다른 구분건물의 소유권의 등기명의인을 대위하여 그 등기를 신청할 수 있다.
  \item 건물이 구분건물인 경우에 그 건물의 등기기록 중 1동 표제부에 기록하는 등기사항에 관한 변경등기는 그 구분건물과 같은 1동의 건물에 속하는 다른 구분건물에 대하여도 변경등기로서의 효력이 있다.
  \item 1동의 건물에 속하는 구분건물 중 일부만에 관하여 소유권보존등기를 신청하는 경우에는 나머지 구분건물의 표시에 관한 등기를 동시에 신청하여야 한다.
\end{choices}
\end{question}
\answer{1}
